\section{Conclusion}
Source-relative recording and ordered tie resolution permit cross-group scheduling when each native flood retains private state. Whole-pair FP32 waves expose a distinct opportunity within each flood, with explicit numerical admission and bucket closure. The input audit establishes three two-layer graphs and the BNN profile quantifies the work induced by coarser scheduling. On these two-layer graphs, the same-rule CPU/CUDA comparison yields construction speedups of 3.31 and 3.18. Ordered STRICT consumption provides a conditional native-map guarantee, while ALT requires independent route qualification.
The 36 primary GEMM and BNN runs establish matched performance. LSTM extends the auxiliary map and backend checks.
